%%% REVISION NOTES (v2, combined Moews + Shrager) %%%
%%%
%%% Per Alex Magoun's editorial guidance:
%%%   - Folded in Moews's IPL-I (LT1) reanimation; authorship is
%%%     alphabetical (see footnote on author block).
%%%   - Cut the IPL-Lisp comparison (old Sec. 2.2 second half and the
%%%     "Reflections: IPL, Lisp, and the Arc of AI" section).
%%%   - Reduced the LLM-debugging material to one short subsection;
%%%     the full narrative is reserved for a possible later submission.
%%%   - Added Sec. "What Executable Archaeology Contributes" to answer
%%%     Alex's question directly.
%%%
%%% TODO (DM): Please check every statement about the IPL-I work for
%%% accuracy; placeholders are marked "TODO (DM)". In particular:
%%%   - results table / which *2 theorems LT1 proves under what limits
%%%   - exact line counts for LT1 vs LT5 listings
%%%   - description of the patches and bug fixes
%%%   - your affiliation and email
%%%   - acknowledgments
%%%
%%% TODO (both): The date of the first machine proof. The draft
%%% previously said machine proofs came from IPL-II runs "late 1956 and
%%% early 1957"; standard accounts (and DM's README, citing Simon's
%%% Models of My Life and Newell & Shaw 1957) put the first JOHNNIAC
%%% proof in August 1956. Text below now says August 1956 -- please
%%% verify against the archival record.
%%%
%%% Carried over from v1 notes, still open:
%%%   - Dedicate to Richard Wexelblat?
%%%   - Mention John Hessler's project at Johns Hopkins (prelim LT work)
%%%   - Dick's wonderful paper (at beginning of history)

\documentclass[journal]{IEEEtran}

\usepackage{cite}
\usepackage{amsmath}
\usepackage{graphicx}
\usepackage{url}
\usepackage{hyperref}
\usepackage{booktabs}
\usepackage{listings}

\lstset{
  basicstyle=\ttfamily\small,
  breaklines=true,
  frame=single,
  xleftmargin=1em,
  framexleftmargin=0.5em
}

\begin{document}

\title{Executable Archaeology: Reanimating the Logic Theorist from its
  IPL-I and IPL-V Sources}

\author{David~Moews and Jeff~Shrager%
\thanks{Authorship is alphabetical; the authors contributed
  independently and equally, as described in the text.}%
\thanks{D. Moews: \texttt{dmoews@fastmail.fm}.} % TODO (DM): affiliation?
\thanks{J. Shrager is with Bennu Climate, Inc.:
  \texttt{jshrager@gmail.com}.}}

\maketitle

\begin{abstract}

The Logic Theorist (LT), created by Allen Newell, J.~C.~Shaw, and
Herbert Simon in 1955--1956, is widely regarded as the first
artificial intelligence program. Its history spans several distinct
incarnations as the underlying Information Processing Language (IPL)
evolved: the original 1956 IPL-I ``pseudocode,'' which was only ever
executed by hand; the IPL-II version that ran on the JOHNNIAC; and
Stefferud's 1963 pedagogical re-coding in IPL-V. Here we describe two
independent reanimations that bracket this history. Moews implemented
an interpreter for the IPL-I abstract machine and, after minor
repairs, executed the original 1956 listing---to our knowledge the
first time the IPL-I Logic Theorist has ever been run by a
machine. Shrager implemented a complete IPL-V interpreter in Common
Lisp and executed the code transcribed directly from Stefferud's 1963
RAND report, which proved 16 of 23 attempted theorems from Chapter~2
of \textit{Principia Mathematica}---results historically consistent
with the original system's behavior. Comparing the two reanimations
makes empirically visible the gap between the 1956 cognitive
specification, which depended on the intelligence of its human
simulators, and the engineering required to make a real machine prove
theorems end-to-end.

\end{abstract}


\begin{IEEEkeywords}
Logic Theorist, IPL-I, IPL-V, artificial intelligence history,
software archaeology, A.~Newell, H.~A.~Simon, J.~C.~Shaw, list
processing, digital preservation
\end{IEEEkeywords}

\section{Introduction}
\label{sec:intro}

In the summer of 1956, a small group of researchers gathered at
Dartmouth College for a workshop on what John McCarthy had dubbed
``artificial intelligence.'' Among them were Allen Newell and Herbert
Simon, who arrived with something none of the other attendees had: a
program---albeit one not yet running on any machine---that could
discover proofs of mathematical theorems.\footnote{It is worth noting
that Art Samuel's checkers-playing program existed a couple of years
before LT, and it too was heuristic. But Samuel's goal was to build a
checkers player, whereas Newell, Simon, and Shaw's explicit goal was
to model human cognition across many domains. Samuel's effort was
brilliant in its own right---he also built arguably the first machine
learning program, and the first system to learn from self-play.}

The Logic Theorist (LT)\cite{LT1956a,LT1956b} would ultimately prove
38 of the first 52 theorems in Chapter~2 of Whitehead and Russell's
\textit{Principia Mathematica}~\cite{whitehead1910principia}, in some
cases finding proofs more elegant than those produced by the human
authors. For Theorem~2.85, LT discovered a shorter, more direct proof
than the one published in \textit{Principia}. Simon showed the new
proof to Bertrand Russell himself, who ``responded with
delight''~\cite[p.~167]{mccorduck2004machines}. Edward Feigenbaum,
then an undergraduate taking a graduate course from Simon, recalls the
moment the project's significance was first announced. As he later
told Pamela McCorduck: ``It was just after Christmas vacation---January
1956---when Herb Simon came into the classroom and said, `Over
Christmas Allen Newell and I invented a thinking machine.' And we all
looked blank''~\cite[p.~138]{mccorduck2004machines}.

Beyond reproducing \textit{Principia} proofs, the Logic Theorist
introduced heuristic search strategies---means--ends analysis and
subgoal decomposition---derived from protocol analyses of human
problem solving. These ideas later became foundational in both
artificial intelligence and cognitive architectures such as the
General Problem Solver (GPS)~\cite{newell1959gps}, Soar, ACT, and
EPAM~\cite{feigenbaum1961epam}. LT thus sits at the very beginning of
the fusion of AI and cognitive psychology.

LT was implemented in the Information Processing Language (IPL), a
language created by the same team---Newell, Shaw, and
Simon---specifically for building cognitive
models~\cite{newell1964ipl}. IPL introduced list manipulation, dynamic
memory allocation, symbols as first-class objects, recursion, and
higher-order functions, and was an explicit, acknowledged ancestor of
Lisp~\cite{mccarthy1960lisp} and of other early list-processing
systems such as Weizenbaum's SLIP. Yet while Lisp has thrived for
nearly seven decades, IPL has been extinct since the mid-1960s.
Although numerous IPL-V implementations existed in the 1960s on
various machines, no maintained implementations appear to be available
in modern environments. The only prior modern attempt known to the
authors is Richman's unpublished IPL-V interpreter, written in
Microsoft BASIC, which he used to run EPAM~IV in the early
1990s~\cite{richman1995epam}.\footnote{Richman wrote to the second
author (personal communication, Sep~26, 2025, quoted with
permission): ``I just interpreted IPL-V in Microsoft Basic using the
IPL-V manual that I found at the CMU library.''}

This paper describes two reanimations of the Logic Theorist,
undertaken independently and only later discovered to be
complementary. Moews implemented an interpreter for IPL-I---the 1956
``pseudocode'' in which LT was first published, and which was never
implemented on any machine---and, after minor repairs to the printed
listing, executed the original 1956 program. Shrager implemented a
complete IPL-V interpreter in Common Lisp and executed the LT code
transcribed directly from Stefferud's 1963 RAND technical
report~\cite{stefferud1963lt}, in which more than 99\% of the running
code is the original code from that report. Following the convention
that the version number tracks the IPL generation, we will call the
1956 IPL-I program \emph{LT1} and the 1963 IPL-V program \emph{LT5}.

This work illustrates a methodological approach that might be called
\emph{executable archaeology}: the study of historically significant
computational systems through faithful reconstruction and execution of
their original code. Early AIs were not merely theoretical proposals
but functioning programs (or, in the case of LT1, hand-executed
procedures) whose detailed behavior cannot be fully understood from
published descriptions alone. Reanimating such systems provides a new
empirical window into the origins of artificial intelligence; we
return to exactly what that window shows in
Section~\ref{sec:contribution}.

In prior work, a team including the second author restored
Weizenbaum's original ELIZA to operation on a reconstructed CTSS
running on an emulated IBM 7094~\cite{lane2025eliza}. Although the
Logic Theorist and ELIZA addressed very different problems---formal
theorem proving versus natural-language conversation---both emerged
from the same symbolic, list-processing programming lineage
originating in IPL. However, whereas the ELIZA work rested on a stack
that emulated vintage hardware and operating systems, the present
reanimations emulate the abstract IPL-I and IPL-V machines.

\section{Historical Background}
\label{sec:history}

\subsection{The Origins and Versions of the Logic Theorist}

The story of the Logic Theorist begins in the early 1950s at the RAND
Corporation. Herbert Simon and Allen Newell, collaborating with
J.~C. (Cliff) Shaw, sought to simulate complex thought by manipulating
symbols. However, the history of the program is characterized by a
series of shifting specifications across different versions of the
Information Processing Language.

The system described in the famous 1956 RAND report
(P-868)~\cite{LT1956a}, and in nearly identical form in the IRE
Transactions that September~\cite{LT1956b}, was a ``paper'' version,
written in what was called IPL-I or the ``Logic Language'' (LL). IPL-I
was never implemented on any machine: it was a pseudocode for an
imagined abstract machine, executed entirely by hand. At the time of
the Dartmouth workshop in the summer of 1956, Newell and Simon noted
that the program was ``essentially specified'' and that hand
simulations suggested it would prove the 60-odd theorems of Chapter~2
of \textit{Principia Mathematica}. These hand simulations famously
employed ``the human computer'': Simon's wife, his children, and
graduate students acting as program
components~\cite[Ch.~13]{simon1996models}.

The first version to run on real hardware was written in IPL-II for
the RAND JOHNNIAC~\cite{gruenberger1968johnniac,newellshaw1957},
producing its first machine proof in August
1956~\cite{simon1996models,newellshaw1957}. The widely cited result
that LT proved 38 of the first 52 theorems in
\textit{Principia}---including the discovery of a more elegant proof
for Theorem 2.85 than the one produced by Whitehead and
Russell---stems from the empirical explorations reported in
1957~\cite{LT1957}.

The IPL-V version used in Shrager's reanimation, transcribed from
Stefferud's 1963 memorandum, represents a third distinct phase. The
program was converted to IPL-V by Fred Tonge and further developed by
Stefferud as a pedagogical re-coding of the original heuristic logic
into the first version of the language made widely available for
public use~\cite{stefferud1963lt}. Thus the two reanimations described
here bracket the program's history: LT1 is the 1956 conceptual model
exactly as first published, while LT5 is the 1963 ``modernized''
production of that model.

%%% TODO (both): confirm/adjust years in the 2026 rows.
\begin{table}[h]
\centering
\caption{Evolution of the Logic Theorist and IPL}
\label{tab:lt_versions}
\begin{tabular}{@{}llll@{}}
\toprule
\textbf{Year} & \textbf{Version} & \textbf{Platform} & \textbf{Historical Milestones} \\ \midrule
1956 & IPL-I (LL) & Manual & Conceptual spec; hand-simulated only. \\
1956--57 & IPL-II & JOHNNIAC & First machine proofs; 38 of 52 theorems. \\
1958 & IPL-IV & IBM 704 & Internal RAND use; transition to IPL-V. \\
1963 & IPL-V & Various & Stefferud's pedagogical LT implementation. \\
2026 & IPL-V & Lisp & Reanimation of LT5 (Stefferud's source). \\
2026 & IPL-I & Python & First machine execution of LT1. \\ \bottomrule
\end{tabular}
\end{table}

\subsection{IPL: The Invention of Symbol- and List-Processing}

To implement LT, Newell, Shaw, and Simon needed a programming language
suited to symbolic reasoning. No such language existed, so they
created the ``Information Processing Language.'' As noted, IPL went
through several iterations. The earliest versions (IPL-I and IPL-II)
powered the 1956 hand simulations and subsequent JOHNNIAC runs;
IPL-V, released in the early 1960s~\cite{newell1964ipl}, became the
definitive public standard.

IPL-V runs on an abstract machine organized around \textit{cells}.
Each cell has an associated symbol. There is no single distinguished
stack; every cell is potentially a stack of symbols that can be pushed
and popped. Symbols are central to IPL-V: they are first-class
objects, and every cell, data structure, routine, and control element
is addressable by a symbol (its name). A set of control cells
(H0--H5) governs execution: H0 serves as the parameter and result
``communication'' cell, H1 serves as both program counter and call
stack, H2 is the free list for dynamic memory allocation, and H3 is
the cycle counter. Ten working cells (W0--W9) provide scratch
storage. Everything in IPL-V---data, routines, even the program
itself---is a list of cells, and thus also a list of symbols. Routines
are invoked by pushing their symbol (name) onto H1. The IPL machine
interprets the operation indicated in that cell, and then continues to
whatever is the indicated next cell. Elegantly, because H1 is both
program counter and call stack, no special machinery is required to
operate the machine aside from basic stack manipulation instructions.

The language's built-in operations are called J-functions (because
their names begin with ``J'': J0, J1, J2, and so on). There are
approximately 150 of these, implementing operations from basic list
manipulation to ``generator'' (iterator) protocols. In the original
IPL-V, many J-functions were themselves written in IPL, although some
were implemented at the level of the underlying physical machine.

IPL-I, the language of the 1956 publication, is recognizably the same
kind of machine seen through an earlier lens: an assembler-like
notation for an abstract list-processing computer, with primitive
operations playing the role later standardized as the J-functions.
But because IPL-I was executed only by hand, its specification could
afford to be looser, leaving to the human executor details that a real
machine implementation would have to nail down. This difference turns
out to be central to what the two reanimations jointly reveal
(Section~\ref{sec:contribution}).

\section{Reanimating LT1: The 1956 IPL-I Program}
\label{sec:lt1}

%%% TODO (DM): This section is drafted from your email, README, and
%%% repo structure. Please correct, expand, and add concrete detail,
%%% especially: the nature of the patches/bug fixes, how the
%%% interpreter handles IPL-I's looseness, results on the *2 theorems,
%%% and anything notable about hand-simulation-era assumptions you
%%% encountered.

The source code for LT1 was published twice in 1956, in RAND paper
P-868~\cite{LT1956a} and in the IRE Transactions on Information
Theory~\cite{LT1956b}; the listings are almost identical. The code is
written in IPL-I, an assembler-like language for an abstract machine
that, as discussed above, was never implemented on non-human
hardware. The listing comprises roughly 400 lines.
%%% TODO (DM): exact count, and which printing you treat as primary.

Moews transcribed the listing from the published report and
implemented an interpreter for the IPL-I abstract machine in Python.
The printed code contains typos and other problems and is not
immediately runnable in its published form; it required minor patches
and bug fixes, each documented in the repository's source
files~\cite{moewsrepo}. After these repairs, the 1956 program runs.
To our knowledge, this is the first time the IPL-I version of the
Logic Theorist has ever been executed by a machine.

One omission in the published program is particularly telling.
Although the IPL-I code is capable of finding individual proof steps,
and of determining when enough steps have been found to constitute a
proof, it contains no logic to print the steps or to assemble them
into an actual proof. The human simulators, of course, needed no such
logic---they could see the steps they had executed and write the proof
down. To obtain proofs from the machine execution, the interpreter
prints each potential proof step as it is found, and a separate
program assembles the proof from the resulting set of steps and checks
it~\cite{moewsrepo}.

A driver program imitates the 1957 ``Empirical
Explorations''~\cite{LT1957} regime by running LT1 successively on
each of the theorems of \textit{Principia} Chapter~2, allowing it to
use, in the proof of each theorem, all axioms and all previous
Chapter~2 theorems, whether or not it managed to prove them. In
addition to substitution and detachment, LT1 employs chaining,
inferring $p \rightarrow r$ from $p \rightarrow q$ and $q \rightarrow
r$.

%%% TODO (DM): Insert a results summary here (how many of the *2
%%% theorems LT1 proves, under what limits, and how this compares both
%%% to the 1957 reported results and to LT5's results in
%%% Table~\ref{tab:results}). A small table parallel to
%%% Table~\ref{tab:results} would be ideal. An example LT1 proof,
%%% parallel to the LT5 trace in Section~\ref{sec:results}, would also
%%% strengthen the comparison.

\section{Reanimating LT5: Stefferud's 1963 IPL-V Program}
\label{sec:lt5}

\subsection{Source Materials and Transcription}
\label{sec:sources}

Unlike the ELIZA reanimation, which began with the dramatic discovery
of original source code printouts in Weizenbaum's archives at
MIT~\cite{lane2025eliza}, the source materials for the Logic Theorist
were already publicly available. Two documents were indispensable for
the LT5 work:

\begin{enumerate}
\item \textbf{Stefferud (1963)}, \textit{The Logic Theory Machine: A
  Model Heuristic Program}, RAND Corporation memorandum
  RM-3731~\cite{stefferud1963lt}. This report contains a complete
  listing of the LT program in IPL-V---nearly 3,000 lines---along with
  detailed descriptions of every aspect of LT's workings, including
  numerous flow charts.

\item \textbf{Newell et al.\ (1964)}, \textit{Information Processing
  Language-V Manual}, 2nd edition~\cite{newell1964ipl}. The definitive
  reference for IPL-V, including the semantics of all J-functions,
  cell structure conventions, and the generator protocol, as well as
  numerous student exercises and sample programs.
\end{enumerate}

Both documents proved to be exceptionally well-written. There were,
however, occasional ambiguities---places where the specification
assumed familiarity with conventions of 1960s computing (such as BCD
character encoding) that are no longer common knowledge.

The LT code was transcribed from Stefferud's report into a spreadsheet
by Shrager and Anthony Hay. Subsequently, Rupert Lane identified
several transcription errors using his ``gridlock''
tool~\cite{lane2024gridlock}, software designed to accurately extract
tabular code from PDF documents. Despite multiple people putting
multiple sets of eyes on the code, a surprising number of small
transcription errors persisted through numerous rounds of review.

The transcribed code was converted into an S-expression format
(\texttt{.liplv} files) suitable for loading by the Lisp-based
interpreter. Conveniently, Lisp comment characters work in these
files, allowing corrections and annotations to be documented inline.

\subsection{A New IPL-V Interpreter in Lisp}
\label{sec:interpreter}

Shrager implemented a new IPL-V interpreter, in Common Lisp, to
execute the Logic Theorist listing. The interpreter implements the
symbolic data structures and execution model described in the IPL-V
manual and supports the primitives required by the program. Although
the earliest versions of the Logic Theorist ran on the RAND JOHNNIAC,
by the time the Stefferud listing was published (1963) IPL-V had
already been implemented on multiple machines. Accordingly, this
reconstruction executes the program within a machine-independent IPL-V
interpreter rather than attempting to reproduce a specific JOHNNIAC
hardware environment.

There were several reasons for emulating the IPL-V abstract machine
rather than the original hardware environment (as was done for ELIZA
on the IBM 7094 and CTSS~\cite{lane2025eliza}). First, the code in
question is IPL-V, and so is not JOHNNIAC-specific. Additionally, no
thorough JOHNNIAC emulator exists, and in any case the JOHNNIAC IPL
code is not known to survive. There does exist an IBM 7094 emulator
that includes an IPL-V interpreter running under a variant of the
IBSYS batch operating system~\cite{ibsys}, and this emulator includes
demo scripts for IPL-V. However, that IPL-V interpreter is a binary
with no source code and no documentation; an attempt to run LT on this
platform failed early in execution with no useful information about
the cause, and with no source code the failure was essentially
undebuggable.

More fundamentally, a central goal of the project was not merely to
get LT to run, but to \textit{elucidate} IPL-V---to develop a deep
understanding of the language, the machine model, and the
infrastructure of the period. Writing a new interpreter was essential
to that goal. IPL-V is defined as an abstract machine specification,
not as a program for a specific computer; it was implemented on many
different machines through the mid-1960s before being supplanted by
Lisp.

The interpreter (\texttt{iplv.lisp}) implements the complete IPL-V
abstract machine as described in the 1964 manual: the cell table, the
control cells H0--H5, the working cells W0--W9, and the push-down
mechanisms. Lists are stored as linked chains of cells, following the
manual's specification. It is worth noting that the push-down (stack)
mechanism is itself implemented as linked lists within the IPL-V
system. Push and pop operations, which appear to be hardware
primitives of the abstract machine, are internally built on the same
list mechanism that underlies everything else in IPL-V, including data
structures, routines, and IPL programs themselves; like Lisp, IPL-V is
lists all the way down.

The core of the interpreter is a function, \texttt{ipl-eval}, that
fetches and executes instructions by executing the list named in the
program counter (H1). The evaluator is re-entrant---a capability
trivially implementable in Lisp---supporting the recursive calls that
IPL-V requires. The J-functions are implemented as Lisp functions.
While the interpreter implements the full IPL-V machine, only the
subset of J-functions required by LT (and a few test programs) is
currently implemented, amounting to several dozen of the approximately
150 defined in the manual.

A complete run of LT5 attempting proofs of 23 theorems is available at
\cite{runoutput}. This takes approximately 574,000 IPL machine cycles
and creates over 576,000 symbols\footnote{Yeah, we were surprised too.
It's really just a coincidence!}. It finishes in a few seconds on a
typical modern laptop. On the JOHNNIAC at RAND, in the mid 1950s, this
likely would have required hours.

\subsection{Debugging Against the Archival Record}
\label{sec:debugging}

%%% Note: the extended account of LLM-assisted debugging has been cut
%%% per editorial guidance and may seed a separate submission. The
%%% Simon's J's episode is retained because it bears directly on the
%%% executable-archaeology argument.

The most difficult phase of debugging the LT5 reanimation began once
the interpreter was stable enough that LT ran without crashing: the
program would prove trivial theorems but fail on simple ones, fall
into runaway recursions, or fail to recognize a completed
proof---anomalies whose causes were far from their symptoms. This
experience has an uncanny resonance with the original developers'
accounts. In an interview conducted by Pamela McCorduck for her
pioneering history of AI, Cliff Shaw described the most daunting phase
of troubleshooting as occurring late in development, when the system
was beyond crashes and produced behavior that appeared superficially
reasonable, leaving the developers in a quandary as to whether the
results were valid consequences of the heuristic logic or actual
bugs, in a memory that had become a churning and impenetrable web of
interconnected lists~\cite{cmuarchiveshawinterview}. Sixty years
later, with the same program, Shrager faced the identical challenge.
In this latter phase, large language models (Google's Gemini and
Anthropic's Claude) were used as debugging assistants, performing the
grinding work of reading thousands of lines of execution traces and
tracking machine state across tens of thousands of cycles; a fuller
account of that collaboration is beyond the scope of this paper.

One episode, however, bears directly on the methodology at issue
here. The final substantive bug was localized to J74, the list-copying
function, yet the Lisp implementation appeared to match the manual's
prose description; the specification itself seemed to support the
buggy behavior. Resolution came not from documentation but from an
artifact: ``Simon's J's,'' a set of original IPL-V punched cards from
1962 preserved by the Computer History
Museum~\cite{bochannek2012simonsjs}, representing the production
implementation of the J-functions at the moment IPL-V was being
standardized. Comparing the 1962 card-deck implementation of J74 with
the modern one confirmed a hypothesized discrepancy, and a one-line
patch made LT run completely correctly. The manual's prose, a careful
modern reading of it, and inference from the modern code had all been
insufficient; the original 1962 source was the ground truth.

\section{What the Logic Theorist Proves}
\label{sec:results}

The Logic Theorist attempts to prove theorems in propositional logic
using the five axioms of \textit{Principia Mathematica}, Chapter~2.
The notation used by the original authors is given in
Table~\ref{tab:notation}.

\begin{table}[h]
\centering
\caption{Logical notation used by the Logic Theorist}
\label{tab:notation}
\begin{tabular}{cl}
\toprule
\textbf{Symbol} & \textbf{Meaning} \\
\midrule
\texttt{I} & Implication ($\rightarrow$) \\
\texttt{V} & Disjunction ($\lor$) \\
\texttt{*} & Conjunction ($\land$) \\
\texttt{-} & Negation ($\lnot$) \\
\texttt{=} & Biconditional ($\leftrightarrow$) \\
\texttt{.=.} & Definitional equality \\
\bottomrule
\end{tabular}
\end{table}

LT employs several proof methods:

\begin{itemize}
\item \textbf{Substitution}---substitute variables in an axiom or
  previously proved theorem to match the goal.
\item \textbf{Detachment} (modus ponens)---from $A \rightarrow B$ and
  $A$, conclude $B$.
\item \textbf{Forward chaining}---apply substitution and detachment
  working forward from known theorems.
\item \textbf{Backward chaining}---work backward from the goal,
  seeking premises that would yield it.
\item \textbf{Sublevel replacement}---replace a subexpression using a
  known equivalence.
\end{itemize}

In LT5, each proof attempt is bounded by a settable effort limit in
terms of IPL machine cycles. (In all examples below this limit is set
to 20,000 cycles.) This is a soft cap: it prevents the system from
\textit{beginning} a new subproblem exploration once exceeded, but
does not interrupt a search in progress.
%%% TODO (DM): describe LT1's effort-limiting (or lack thereof) for
%%% parallelism.

Table~\ref{tab:axioms} shows the axioms and definitions that are
initially provided to LT5 for the runs described in the present
paper. Importantly, as LT proves a given theorem, that theorem is
treated as an axiom and can be used to support subsequent proofs.
Thus, LT can be said, in a sense, to learn!

\begin{table}[h]
\centering
\caption{Axioms and definitions}
\label{tab:axioms}
\begin{tabular}{ll}
\toprule
\textbf{Name} & \textbf{Axiom/Def}\\
\midrule
1.01  & \texttt{((PIQ).=.(-PVQ))}\\
2.33  & \texttt{((PVQVR).=.((PVQ)VR))}\\
3.01  & \texttt{((P*Q).=.-(-(PV-Q))}\\
4.01  & \texttt{((P=Q).=.((PIQ)*(QIP)))}\\
1.2   & \texttt{((AVA)IA)}\\
1.3   & \texttt{(BI(AVB))}\\
1.4   & \texttt{((AVB)I(BVA))}\\
1.5   & \texttt{((AV(BVC))I(BV(AVC)))}\\
1.6   & \texttt{((BIC)I((AVB)I(AVC)))}\\
\bottomrule
\end{tabular}
\end{table}

Table~\ref{tab:results} shows the results of the reanimated LT5 on the
23 theorems attempted: 16 are proved, and the remaining 7 are
not---results that are historically consistent with the original LT's
known behavior, given an effort limit of 20,000 IPL machine cycles.

\begin{table}[h]
\centering
\caption{Theorems attempted (LT5)}
\label{tab:results}
\begin{tabular}{lll}
\toprule
\textbf{Name} & \textbf{Formula} & \textbf{Result} \\
\midrule
2.01 & \texttt{(PI-P)I-P} & Proved \\
2.02 & \texttt{QI(PIQ)} & Proved \\
2.04 & \texttt{(PI(QIR))I(QI(PIR))} & Proved \\
2.05 & \texttt{(QIR)I((PIQ)I(PIR))} & Proved \\
2.06 & \texttt{(PIQ)I((QIR)I(PIR))} & Proved \\
2.07 & \texttt{PI(PVP)} & Proved \\
2.08 & \texttt{PIP} & Proved \\
2.10 & \texttt{-PVP} & Proved \\
2.11 & \texttt{PV-P} & Proved \\
2.12 & \texttt{PI--P} & Proved \\
2.13 & \texttt{PV---P} & Proved \\
2.14 & \texttt{--PIP} & Not proved \\
2.15 & \texttt{(-PIQ)I(-QIP)} & Not proved \\
2.20 & \texttt{PI(PVQ)} & Proved \\
2.21 & \texttt{-PI(PIQ)} & Not proved \\
2.24 & \texttt{PI(-PVQ)} & Proved \\
3.13 & \texttt{(-(P*Q))I(-PV-Q)} & Not proved \\
3.14 & \texttt{(-PV-Q)I(-(P*Q))} & Not proved \\
3.24 & \texttt{-(P*-P)} & Not proved \\
4.13 & \texttt{P=--P} & Not proved \\
4.20 & \texttt{P=P} & Proved \\
4.24 & \texttt{P=(P*P)} & Proved \\
4.25 & \texttt{P=(PVP)} & Proved \\
\bottomrule
\end{tabular}
\end{table}

A representative proof trace is shown below. Here LT5 proves Theorem
2.01, \texttt{(PI-P)I-P}:

\begin{lstlisting}
2.01    (PI-P)I-P

PROOF FOUND.

    GIVEN           *1.2  (AVA)IA
    SUBSTITUTION    .0    (PI-P)IP
    SUBLEVEL REPL   2.01  (PI-P)I-P
    Q.E.D.

EFFORT      LIMIT 20000   ACTUAL 5579
SUBPROBLEMS LIMIT 50      ACTUAL 1
SUBSTITUTIONS LIMIT 50    ACTUAL 2
\end{lstlisting}

The proof begins from Axiom~*1.2, \texttt{(AVA)IA}, substitutes
\texttt{-P} for \texttt{A} to obtain \texttt{(-PV-P)I-P}, then applies
sublevel replacement (using the fact that \texttt{P} implies
\texttt{PVP}) to transform the antecedent, yielding the desired
theorem. The Appendix depicts a more complex proof (4.25), and the
original output, including both of these proofs, is at
\cite{runoutput}.

%%% TODO (DM): add a parallel LT1 trace or results note here or in
%%% Section~\ref{sec:lt1}, so the reader can compare the two systems'
%%% behavior on the same theorem(s).

\section{What Executable Archaeology Contributes}
\label{sec:contribution}

%%% This section is new, written in response to Alex's question:
%%% "What exactly does it contribute to understanding the origins of
%%% AI in this case?"

It is reasonable to ask what is learned by running this old code that
could not be learned by reading it, or by reading the abundant
secondary literature on the Logic Theorist. The present case suggests
several answers.

\textbf{It makes the gap between theory and implementation
measurable.} LT1 and LT5 are, in intent, the same program; in fact
they differ by nearly an order of magnitude in size---roughly 400
lines of IPL-I against nearly 3,000 lines of
IPL-V.\footnote{%%% TODO (DM/JS): replace with exact counts and state
the counting convention (instructions? card images?).} The 1956
program is something close to a direct expression of Newell and
Simon's hypothesized cognitive operators for human theorem proving: a
high-level specification of a theory of cognition. It could afford to
be that, because its ``hardware'' was human---intelligent, flexible,
and able to silently fill in whatever the specification left
unsaid. LT5 is what the same design became when Shaw and his
successors had to make a real, literal-minded machine do the whole
thing end-to-end. Reanimating both versions converts a familiar
qualitative observation about early AI---that getting from idea to
running program was most of the work---into something that can be
measured, localized, and studied line by line.

\textbf{It reveals exactly what the human computer was silently
doing.} The most striking single finding of the LT1 reanimation is an
absence: the published 1956 program can find proof steps and can
recognize that a proof exists, but contains no machinery to assemble
or output the proof itself. No reader of the report appears to have
remarked on this, and no reading of the secondary literature would
surface it; it only becomes visible when a machine, rather than one of
Simon's students, executes the code and the proof fails to
appear. The omission is itself historical evidence: it marks
precisely a boundary of what Newell and Simon considered the
\emph{theory} (finding the proof) versus mere clerical work
(writing it down)---work the human simulators performed without
anyone needing to specify it.

\textbf{It converts historical claims into checkable ones.} The 1956
report claimed, on the basis of hand simulation, that the program
would prove the theorems of \textit{Principia} Chapter~2; the 1957
paper reported the famous 38-of-52 result for the JOHNNIAC version.
Until now, both claims rested on testimony and on archival output
listings. The reanimations allow such claims to be tested directly:
LT5 proves 16 of 23 attempted theorems under a fixed effort limit, a
result consistent with the original reports, and LT1's behavior on
the same theorem set can now be examined for the first
time.
%%% TODO (DM): one sentence on how LT1's results bear on the 1956
%%% "hand simulations suggest it will prove them" claim.
Moreover, both systems are now experimental instruments: effort
limits, axiom orderings, and theorem sequences can be varied, and the
consequences observed on the genuine historical artifact rather than
on a modern paraphrase of it.

\textbf{It shows that some semantics survive only in artifacts, not
in documentation.} The J74 episode (Section~\ref{sec:debugging}) is a
small but pointed demonstration: the IPL-V manual's prose description
of a core primitive was consistent with an incorrect implementation,
and the discrepancy could be resolved only by consulting the original
1962 punched-card source preserved by the Computer History Museum.
For the history of computing this is a methodological moral:
specifications and manuals, however well written, systematically
under-determine the systems they describe, and executable
reconstruction is one of the few ways to discover where.

\textbf{It corrects the published record.} Both source listings turn
out to contain errors---typos and inconsistencies in the 1956 IPL-I
listing that prevent it from running as printed, and subtle
transcription hazards in working from the 1963 report---now documented
in the respective repositories~\cite{moewsrepo,runoutput}. These
repairs, and the requirement that every repair be justified and
recorded, are the ``archaeology'' in executable archaeology.

Finally, the two reanimations were undertaken independently, by
authors unknown to each other, and converged on complementary
reconstructions of adjacent stages of the same program. This
independence functions as a kind of replication: agreement between
the two efforts on the program's structure and behavior is evidence
that the reconstructions reflect the historical system rather than
either reconstructor's assumptions.

\section{Related Work}
\label{sec:related}

The reanimation of historical software systems has gained momentum in
recent years. Most directly related is the ELIZA reanimation by Lane
et al.~\cite{lane2025eliza}, in which a team including the second
author restored Weizenbaum's original ELIZA to operation on a
reconstructed CTSS running on an emulated IBM~7094. That project
reconstructed the complete original hardware and operating system
stack; the present work instead emulates abstract machine
specifications in modern languages.

Reimplementing the Logic Theorist's \textit{algorithms} in a modern
language is a relatively straightforward exercise---essentially an
introductory AI course project, and one that has surely been done many
times. A 1968 paper, for example, describes an implementation of LT in
LISP~1.5~\cite{millstein1968logic}. The present work is a
fundamentally different undertaking: rather than rewriting LT's logic
in a modern language, it constructs faithful emulators of the IPL-I
and IPL-V abstract machines and runs the original code, transcribed
directly from the 1956 and 1963 reports. The distinction is
significant: a Lisp reimplementation demonstrates that LT's algorithms
can be expressed in Lisp (which is unsurprising), while the present
work demonstrates that the original IPL code, as written and published
by Newell, Simon, and Shaw's team, functions as described---and,
moreover, enables experiments on the real, working systems of machine
and program.

\section{Next Steps}
\label{sec:next}

The IPL-V interpreter developed for this project is not specific to
the Logic Theorist. It implements the general IPL-V abstract machine
and can, in principle, run any IPL-V program. Work is in progress on
a corpus of early IPL-V AI programs. The most significant targets are
the General Problem Solver (GPS)~\cite{newell1959gps} and EPAM
(Elementary Perceiver and Memorizer)~\cite{feigenbaum1961epam}. EPAM
source code is available in the Feigenbaum archives at Stanford, and
archival work to locate complete source code for GPS is in progress.
Now that proven IPL interpreters exist, running these programs should
require substantially less effort than the initial construction of
the interpreters themselves.

The IPL-I interpreter, repaired LT1 source, and proof-assembly tools
are open sourced at \url{https://github.com/dmoews/logic-theorist}.
The IPL-V interpreter, the transcribed LT5 code, and complete outputs
as described in this paper are open sourced at
\url{https://github.com/jeffshrager/IPL-V}.

\section*{A Personal Coda}

One of us (Shrager) studied under both Newell and Simon at Carnegie
Mellon University in the 1980s, where his work focused on modeling the
cognition of scientists and engineers. IPL and LT were long since
history by that time, and neither came up in classes or
conversation. But the intellectual values Newell and Simon
instilled---rigor about mechanism, respect for the complexity of
cognition, and the conviction that thinking could be understood
computationally---shaped the careers of their students. There is a
certain closure in now studying Newell, Simon, and Shaw \textit{as}
scientist-engineers: reverse-engineering their design decisions,
reconstructing their reasoning, and trying to understand how they
invented AI and cognitive science. The subject of that training and
its object have folded into one another.

\section*{Acknowledgements}

%%% TODO (DM): please add your acknowledgements.

The following people (listed in no particular order) contributed to
the LT5 effort in ways ranging from code transcription to weekly
encouragement: Anthony Hay, Arthur Schwarz, Leigh Klotz, David
M.\ Berry, Rupert Lane, Amy Majczyk, Emily Davis, Ethan Ableman, and
Paul McJones. Shrager is particularly grateful to Hay, Schwarz, Klotz,
and Berry, who put up with weekly progress updates throughout the
project and responded with encouragement that was always helpful and
sometimes technically substantive. Rupert Lane's gridlock software
caught several transcription errors that had survived multiple rounds
of human review.

We thank the Computer History Museum for preserving and making
available ``Simon's J's,'' the original 1962 IPL-V J-function card
deck, which proved essential during debugging. We also thank Alexander
Magoun and David Walden for editorial guidance, including the
suggestion to combine our independent efforts into this single
account.

Finally, we acknowledge Allen Newell, Herbert Simon, and Cliff Shaw,
whose work this paper merely re-runs.

{\footnotesize Portions of the LT5 debugging effort, and portions of
the drafting of this paper, were carried out with the aid of various
LLMs, but all concepts and final editorial decisions were the
authors' alone.}

\begin{thebibliography}{99}

\bibitem{LT1956a} A. Newell and H. A. Simon, \textit{The Logic Theory
  Machine -- A Complex Information Processing System}, RAND Corp.,
  Santa Monica, CA, USA, Rep. P-868, 1956.

\bibitem{LT1956b} A. Newell and H. A. Simon, ``The logic theory
  machine--A complex information processing system,'' \textit{IRE
  Trans. Inf. Theory}, vol. 2, no. 3, pp. 61--79, Sep. 1956.

\bibitem{LT1957} A. Newell, J. C. Shaw, and H. A. Simon, ``Empirical
  explorations with the logic theory machine: A case study in
  heuristics,'' in \textit{Proc. Western Joint Comput. Conf.}, 1957,
  pp. 218--230.

\bibitem{newellshaw1957} A. Newell and J. C. Shaw, ``Programming the
  logic theory machine,'' in \textit{Proc. Western Joint
  Comput. Conf.}, 1957, pp. 230--240.

\bibitem{whitehead1910principia} A. N. Whitehead and B. Russell,
  \textit{Principia Mathematica}. Cambridge, UK: Cambridge
  Univ. Press, 1910--1913.

\bibitem{newell1964ipl} A. Newell, et al. \textit{Information
  Processing Language-V Manual}, 2nd ed. Englewood Cliffs, NJ, USA:
  Prentice-Hall, 1964.

\bibitem{mccarthy1960lisp} J. McCarthy, ``Recursive functions of
  symbolic expressions and their computation by machine, Part I,''
  \textit{Commun. ACM}, vol. 3, no. 4, pp. 184--195, Apr. 1960.

\bibitem{simon1996models} H. A. Simon, \textit{Models of My Life}.
  Cambridge, MA, USA: MIT Press, 1996.

\bibitem{cmuarchiveshawinterview}
  Pamela McCorduck Collection, Carnegie Mellon University Archives.

\bibitem{ibsys} R. Storey, ``B7094'' [Online]. Available:
  \url{https://github.com/Bertoid1311/B7094} (Accessed: Mar. 11,
  2026).

\bibitem{stefferud1963lt} E. Stefferud, ``The Logic Theory Machine: A
  Model Heuristic Program,'' RAND Corp., Santa Monica, CA, USA,
  Rep. RM-3731-CC, 1963.

\bibitem{lane2025eliza} R. Lane, A. Hay, A. Schwarz, D. M. Berry, and
  J. Shrager, ``ELIZA Reanimated: Restoring the mother of all chatbots
  to one of the world's first time-sharing systems,'' \textit{IEEE
    Ann. Hist. Comput.}, vol. 47, no. 1, pp. 8--23, 2025.

\bibitem{gruenberger1968johnniac} F. Gruenberger, ``The History of the
  JOHNNIAC,'' RAND Corp., Santa Monica, CA, USA, Rep. RM-5654-PR,
  1968.

\bibitem{mccorduck2004machines} P. McCorduck, \textit{Machines Who
  Think}, 2nd ed. Natick, MA, USA: A. K. Peters, 2004.

\bibitem{bochannek2012simonsjs} A. Bochannek, ``Simon's J's,''
  Computer History Museum Blog, Sep. 26, 2012. [Online]. Available:
  \url{https://computerhistory.org/blog/simons-js/} (Accessed:
  Mar. 11, 2026).

\bibitem{lane2024gridlock} R. Lane, ``Gridlock: Accurate code
  extraction from PDFs,'' 2024. [Online]. Available:
  \url{https://github.com/rupertl/gridlock}

\bibitem{millstein1968logic} J. S. Millstein, ``The Logic Theorist in
  LISP,'' \textit{Comput. J.}, vol. 11, no. 1, pp. 39--44, Feb. 1968.

\bibitem{newell1959gps} A. Newell, J. C. Shaw, and H. A. Simon,
  ``Report on a general problem-solving program,'' in
  \textit{Proc. Int. Conf. Inform. Process.}, 1959, pp. 256--264.

\bibitem{feigenbaum1961epam} E. A. Feigenbaum, ``The simulation of
  verbal learning behavior,'' in \textit{Proc. Western Joint
    Comput. Conf.}, 1961, pp. 121--132.

\bibitem{richman1995epam} H. B. Richman, J. J. Staszewski, and
  H. A. Simon, ``Simulation of expert memory using EPAM IV,''
  \textit{Psychol. Rev.}, vol. 102, no. 2, pp. 305--330, 1995.

\bibitem{moewsrepo} D. Moews, ``Recreation of the 1956 IPL-I version
  of the Logic Theorist theorem prover,'' 2026. [Online]. Available:
  \url{https://github.com/dmoews/logic-theorist} (Accessed: Jun. 9,
  2026).

\bibitem{runoutput} J. Shrager, Output from the reanimated logic
  theorist, [Online]. Available:
  \url{https://github.com/jeffshrager/IPL-V/blob/master/major_results/20260309_CompleteAndCorrect.drb}
  (Accessed: Mar. 11, 2026).

\end{thebibliography}

\clearpage % Ensure all previous material is printed
\onecolumn % Switch to single column mode for the appendix

\appendix

\section{Example Complex Logic Theorist Proof}

To illustrate a complex proof generated by the reconstructed Logic
Theorist (LT5), the following derivation corresponds to theorem *4.25
of \textit{Principia Mathematica}:

\[ P = (P \lor P) \]

This states the idempotence of disjunction. In the LT representation,
the equivalence operator is defined using implication and conjunction:

\[ (P = Q) \equiv ((P \rightarrow Q) \land (Q \rightarrow P)) \]

The Logic Theorist proof trace proceeds as follows:

\begin{enumerate}

\item \textbf{Expand the definition of equivalence:}
LT begins by reducing the goal to a conjunction of two implications:
\begin{quote}
\begin{verbatim}
4.25    P=(PVP)
.0   (PI(PVP))*((PVP)IP)  4.01  ,0 REPLACEMENT
\end{verbatim}
\end{quote}
\textit{Logic form:}
\begin{itemize}
    \item[] $P = (P \lor P) \Rightarrow (P \rightarrow (P \lor P)) \land ((P \lor P) \rightarrow P)$ 
\end{itemize}

\item \textbf{Identify the first implication:}
The system identifies a forward-chaining path to establish the first half of the conjunction:
\begin{quote}
\begin{verbatim}
PI(PVP)    4.20  ,0 FORWARD CHAINING
\end{verbatim}
\end{quote}
\textit{Logic form:} 
\begin{itemize}
    \item[] $P \rightarrow (P \lor P)$
\end{itemize}

\item \textbf{Find the proof via substitution:}

LT utilizes Theorem *2.20 from its memory as the basis for
substitution: \textit{Theorem 2.20 is not given as an axiom, but is a
previously proved theorem that it remembered and applied here; See
comment below.}

\begin{quote}
\begin{verbatim}
GIVEN      2.20      AI(AVB)
SUBSTITUTION  .0    PI(PVP)
\end{verbatim}
\end{quote}
\textit{Logic form:} 
\begin{itemize}
    \item[] The general theorem $A \rightarrow (A \lor B)$ is retrieved
    \item[] Substituting $A := P$ and $B := P$ yields $P \rightarrow (P \lor P)$
\end{itemize}

\item \textbf{Complete the equivalence:}
LT completes the chain by linking the proved implications back to the original goal:
\begin{quote}
\begin{verbatim}
GIVEN       4.20      A=A
FORWARD CHAINING    4.25      P=(PVP)
Q.E.D.0
\end{verbatim}
\end{quote}
\textit{Logic form:} 
\begin{itemize}
    \item[] Because both $P \rightarrow (P \lor P)$ and $(P \lor P) \rightarrow P$ are established, the equivalence $P = (P \lor P)$ holds.
\end{itemize}

\end{enumerate}

\textbf{Performance Metrics:}
\begin{quote}
\begin{verbatim}
EFFORT              LIMIT 20000         ACTUAL 8727
SUBPROBLEMS         LIMIT 50            ACTUAL 2
SUBSTITUTIONS       LIMIT 50            ACTUAL 3
\end{verbatim}
\end{quote}

\medskip

\noindent The successful execution of complex derivations such as this
one, which use previously proved theorems, demonstrates the Logic
Theorist's capacity for a primitive form of machine learning: by
treating previously proved theorems, such as 2.20, as established
``givens'' for subsequent search paths, the system bootstraps its own
knowledge base, mirroring the cumulative nature of algebraic
reasoning.

%%% TODO (DM): consider adding a second appendix with an LT1 run on
%%% the same theorem, showing the raw step output and the assembled
%%% proof from analyze_output.py.

\end{document}
